% RMP 2011.12127 fig. 2 acceptance: an MPS (a) and a PEPS (b) with their
% marginals rho_kept = Tr_traced |psi><psi|.
%   (a) top:    the MPS chain |psi> (5 sites, physical legs);
%   (a) bottom: the ket-bra double layer with the RIGHTMOST columns (4, 5)
%               traced out -- kept columns keep open legs (up for ket,
%               down for bra), traced columns close by wrap loops;
%   (b) top:    a 4x4 oblique PEPS sheet, physical=up;
%   (b) bottom: the tenkzplanes double layer with the rightmost column's
%               sites traced.
\documentclass[varwidth=20cm, border=6pt]{standalone}
\usepackage{amsmath, amssymb}
\usepackage{tenkz}
\begin{document}

\textbf{(a) MPS and its marginal}\par\medskip
\begin{tenkz}[physical=down]
  \tn{} & \tn{} & \tn{} & \tn{} & \tn{}
\end{tenkz}
\par\medskip
% rho_{1..3} = Tr_{4,5} |psi><psi|: all-open ket-over-bra window, columns
% 4 and 5 closed per site by the wrap loop (Tr_s)
\begin{tenkz}[rows={ket:nopair, bra}, trace={(physical, 4-5)}]
  \tn{} & \tn{} & \tn{} & \tn{} & \tn{} \\
  \tn{} & \tn{} & \tn{} & \tn{} & \tn{}
\end{tenkz}

\bigskip
\textbf{(b) PEPS and its marginal}\par\medskip
\begin{tenkzlattice}[rows=4, cols=4, frame=oblique, physical=up]
  % The marked segment is one representative virtual contraction between A
  % tensors; all generated bonds and open physical legs keep their roles.
  \tnedge[role=marked]{(2,2)-(2,3)}
\end{tenkzlattice}
\par\medskip
% rho_{1:3} = Tr_{col 4} |psi><psi| acts on the 12 kept sites.  Columns 1--3
% retain one ket and one bra physical
% index at each site, while the four physical pairs in column 4 are traced.
% Ink-to-index correspondence: black vertices on the upper and lower sheets
% are the ket tensors A and conjugate bra tensors; horizontal and diagonal lines
% contract virtual indices within each sheet.  The twelve rising stubs in
% columns 1--3 are the free bra physical indices and the twelve falling stubs
% are the free ket physical indices.  Each wrap loop in column 4 joins the
% ket and bra physical indices at one site, and no other ket--bra segment is
% present.  The lattice frame seals the exterior virtual boundary: suppressed
% outer virtual stubs are fixed boundary contractions, not physical traces.
\begin{tenkzplanes}[rows=4, cols=4, sheets={ket, bra},
                    open={(1-4,1-3)}, trace={(1-4,4)}]
\end{tenkzplanes}

\end{document}
